\paragraph{若干（预）\trEquivariant{}\trBrane{}。}\label{sub:pre-equivariant-branes}
现在我们来回顾\cref{sub:double-branched-cover}中引入的 $\widetilde{E}$ 的那两类\trLagrange\trSubmanifold，目的是把它们变成\trEquivariant\trFukayaCategory{}$\scr{F}(\widetilde{E})^{\mathrm{naive}}$ 的\trObject。回忆一下，为了定义那个\trCategory，我们曾不得不对 $\ZZ/2$-作用施加一些技术性条件，见\cref{thm:fixed-point-set-assumption}。在当前情形下，这些条件显然是满足的：$\iota_{\widetilde{E}}$ 的\trFixedPoint{}集是 $\{y=0\}\cong M$，而且它的\trNormalBundle{}有一个自然的\trTrivialization，由 $D\tilde{\pi}|M\colon T\widetilde{E}|M\to T\widetilde{S}_0=\CC$ 给出。
\begin{itemize}
\item \emph{(U) 型。}设 $L^{\#}=(L,\alpha^{\#},P^{\#})$ 是 $E$ 中的一个\trExactLagrangianBrane，包含在 $E^{\mathrm{shrink}}$ 的内部之中。设 $\widetilde{L}$ 是 $L$ 在 $\widetilde{E}$ 中的\trPreimage。由 \eqref{eq:tilde-e-volume-forms}，$L$ 与 $\widetilde{L}$ 的\trSquaredPhase{}映射之间的关系为
\begin{equation*}
\alpha_{\widetilde{E}}\bigl(T\widetilde{L}_{(y,x)}\bigr)=\frac{|y|^2}{y^2}\,\alpha_E(TL_x).
\end{equation*}
回忆 $\widetilde{L}=\widetilde{L}_+\cup\widetilde{L}_-$，其中\trConnectedComponent{}$\widetilde{L}_{\pm}$ 落在子集 $\pm\mathrm{im}(y)>0$ 之中。由于 $\mathrm{im}(y)$ 永不为零，可以取幅角函数 $\arg\colon\CC\setminus\RR^+\to(0;2\pi)$，并如下定义 $\widetilde{L}$ 的一个\trGrading：
\begin{equation}\label{eq:tilde-l-grading}
\tilde{\alpha}^{\#}(y,x)=\alpha^{\#}(x)-\tfrac{1}{2\pi}\,\arg(y^2).
\end{equation}
限制之后便得到 $\widetilde{L}_{\pm}$ 的\trGrading{}$\tilde{\alpha}^{\#}_{\pm}$。至于\trPinStructure，我们的约定是：给 $\widetilde{L}_+$ 配上已有的 $P^{\#}$ 之\trPullback，而给 $\widetilde{L}_-$ 配上 $P^{\#}\otimes\lambda^{\mathrm{top}}(TL)$ 之\trPullback。把这两者记作 $\widetilde{P}^{\#}_{\pm}\to\widetilde{L}_{\pm}$，并把它们的并记作 $\widetilde{P}^{\#}\to\widetilde{L}$。

有了这些，容易把 $\widetilde{L}_+$ 做成\cref{sub:equivariant-branes}意义下的一个预\trEquivariant{}\trBrane：主分量是 $\widetilde{L}^{\#}_+$，\trShadow{}是 $\widetilde{L}^{\#}_-$，\trIsomorphism{}$\widetilde{f}_-\colon\widetilde{P}^{\#}_-\to\iota_M^*\bigl(\widetilde{P}^{\#}_+\otimes\lambda^{\mathrm{top}}(T\widetilde{L}_+)\bigr)$ 取为恒同，再调整 $\widetilde{f}_+$ 使 \eqref{eq:equivariant-shadow-composition} 恰好成立。把 $+$ 与 $-$ 交换后施行同样的做法，便得到 $\widetilde{L}_-$ 的一个预\trEquivariant{}\trBrane{}结构；取二者的\trDisjointUnion，$\widetilde{L}$ 就成为一个\trEquivariant{}\trBrane。
\end{itemize}

现在我们已经能得出一些初步的推论。设 $L^{\#}_k=(L_k,\alpha^{\#}_k,P^{\#}_k)$（$k=0,1$）是 $E$ 中的两个上述类型的膜。在为这一对选取\trFloerDatum{}$(H,J)$ 时，我们可以要求：在 $E^{\mathrm{shrink}}$ 之外，$H$ 为零，而 $J$ 与 $I_E$ 相一致。这并不妨碍\trModuliSpace{}$\scr{M}_Z(y_0,y_1)$ 的\trRegularity，因为在 $L_0\cup L_1$ 的一个邻域内选取仍然是自由的。对于诸\trPreimage{}$(\widetilde{L}_{0,+},\widetilde{L}_{1,+})$，取把 $(H,J)$ \trLift{}上来所得的\trFloerDatum{}$(\widetilde{H},\widetilde{J})$；这是不成问题的，因为我们的\trFloerDatum{}在分歧轨迹附近是平凡的。每个 $\tilde{y}_+\in\scr{C}(\widetilde{L}_{0,+},\widetilde{L}_{1,+})$ 都落在 $\widetilde{E}^{\mathrm{shrink}}_+$ 之内，并\trProjection{}为某个 $y\in\scr{C}(L_0,L_1)$，这在两个集合之间建立了一个\trBijective。这个对应保持绝对\trIndex，并且（在\trCanonical{}\trIsomorphism{}的意义下）保持\trOrientation{}空间：
\begin{equation}\label{eq:lift-index-orientation}
i(\tilde{y}_+)=i(y),\qquad o(\tilde{y}_+)\cong o(y).
\end{equation}
有一个小的疑点，可以按如下方式处理：虽然 $\eta^2_{\widetilde{E}}|\widetilde{E}^{\mathrm{shrink}}_+$ 并不等于 $\eta^2_E$ 的\trPullback，但通过选取 $\CC\setminus\RR^+$ 到点 $-1$ 上的一个\trDeformationRetract，并把后者应用于 \eqref{eq:tilde-e-volume-forms} 中的项 $|y|^2/y^2$，二者就可以（相差一个整体正负号地）互相\trDeformation。这种二次复\trVolume{}形式之间的\trHomotopy{}诱导出诸\trGrading{}的一个\trHomotopy，而后者事实上把 $\widetilde{L}_k$ 的\trGrading{}$\tilde{\alpha}^{\#}_{k,+}$ \trDeformation{}成 $\alpha^{\#}_k-\tfrac{1}{2}$ 的\trPullback。\trPinStructure{}反正照旧相合，因此 \eqref{eq:lift-index-orientation} 便容易地得到了。对每个 $u\in\scr{M}_Z(y_0,y_1)$，映射 $v=\pi\circ u\colon Z\to S$ 在 $\Omega=v^{-1}(S\setminus S^{\mathrm{shrink}})$ 上是\trHolomorphic{}的，这意味着\trMaximumPrinciple{}适用于 $\overline{\Omega}$ 上的 $|v|$。由于 $\Omega$ 相对紧并且不与 $\partial Z$ 相交，结论是 $\Omega=\varnothing$，于是 $u$ 照旧停留在 $E^{\mathrm{shrink}}_+$ 之内\footnote{译注：原书此处印作 $E^{\mathrm{shrink}}_+$；疑应为 $E^{\mathrm{shrink}}$（$E^{\mathrm{shrink}}$ 本身并无正负号之分，下标当系承上 $\widetilde{E}^{\mathrm{shrink}}_+$ 而误植）；此处照原书排印，存疑备案。——译者注}。这表明，\trProjection{}在上述这些 $u$ 与诸元素 $\tilde{u}_+\in\scr{M}_Z(\tilde{y}_{0,+},\tilde{y}_{1,+})$ 之间建立了一个\trBijective。与之相伴的是相应\trLinearization{}算子的一个等同，故 $(\widetilde{H},\widetilde{J})$ 同样是\trRegular{}的。最后，这一等同与用来定义对\trFloer\trDifferential{} $\partial$ 之贡献 \eqref{eq:relative-orientation-cu} 的那些\trGluing{}步骤是\trCompatible{}的。这就定义了诸\trFloer\trCochainComplex{}的一个\trIsomorphism。一个平行的论证适用于诸\trPreimage{}的 $-$ 部分；唯一的差别在于\trPinStructure{}被 $\lambda^{\mathrm{top}}(TL_k)\cong\xi_E|L_k\cong\xi_{\widetilde{E}}|\widetilde{L}_{k,-}$ 所扭转，但由于这一扭转同等地作用于所涉的两个\trLagrange\trSubmanifold，\trFloerComplex{}并不受影响。总结起来，我们有\trCanonical{}\trIsomorphism
\begin{equation}\label{eq:tilde-cf-decomposition}
\begin{aligned}
CF^*(\widetilde{L}_0,\widetilde{L}_1)&\cong CF^*(\widetilde{L}_{0,+},\widetilde{L}_{1,+})\oplus CF^*(\widetilde{L}_{0,-},\widetilde{L}_{1,-})\\
&\cong CF^*(L_0,L_1)\otimes_\KK\KK[+,-],
\end{aligned}
\end{equation}
其中 $\KK[+,-]$ 表示带有两个\trGenerator{}$+$ 与 $-$ 的二维\trVectorSpace。鉴于诸 $\widetilde{L}^{\#}_k$ 的\trEquivariant{}性质，\trFloerComplex{}带有一个\trCanonical{}\trInvolution；相对于 \eqref{eq:tilde-cf-decomposition}，这正是交换 $+$ 与 $-$ 两个\trGenerator{}的那个映射。更一般地，给定由同类\trLagrange\trBrane{}组成的任一 $(d{+}1)$ 元组，可以选取\trPerturbationDatum，使之具有与上面考虑的\trFloerDatum{}同样的性质。这时，受\trPerturbation{}的\trHolomorphic{} $(d{+}1)$ 边形也将局限于 $E^{\mathrm{shrink}}$ 之内，这意味着 \eqref{eq:tilde-cf-decomposition} 与\trAinfty-结构\trCompatible。
\begin{remark}\label{thm:square-roots-natural-orientation}
假设我们所有的二次复\trVolume{}形式都容许连续的平方根，从而诸膜都自然地带有\trOrientation；见\cref{thm:grading-relative-orientation}。更明确地说，在某一点处\trTangentSpace{}$T\widetilde{L}_{(y,x)}$ 的一组基 $(v_1,\ldots,v_{n+1})$ 是正\trOrientation{}的，当且仅当
\begin{equation*}
\exp\bigl(\pi i\alpha^{\#}(y,x)\bigr)=\frac{\eta_{\widetilde{E}}(v_1\wedge\cdots\wedge v_{n+1})}{\bigl|\eta_{\widetilde{E}}(v_1\wedge\cdots\wedge v_{n+1})\bigr|}.
\end{equation*}
如果我们从 $(y,x)$ 过渡到 $(-y,x)$，$\eta_{\widetilde{E}}=\beta_{\widetilde{E}}^*\eta_E/2y$ 的正负号便会改变，而 $\alpha^{\#}$ 按定义保持不变\footnote{译注：原书两处均印作 $\alpha^{\#}$；疑应为 $\tilde{\alpha}^{\#}$，即 \eqref{eq:tilde-l-grading} 所定义之 $\widetilde{L}$ 的\trGrading{}（按该式，$\arg\bigl((-y)^2\bigr)=\arg(y^2)$，正与下文“按定义保持不变”相合）；此处照原书排印，存疑备案。——译者注}。于是，两个\trConnectedComponent{}$\widetilde{L}_+$ 与 $\widetilde{L}_-$ 带有相反的\trOrientation。
\end{remark}
\begin{itemize}
\item \emph{(B) 型。}设 $\gamma\colon[0;1]\to D$ 是一条以 $\gamma(0)=*$ 为起点的\trVanishingPath。我们要求 $\gamma'(0)\in\CC\setminus\RR^+$，且 $\gamma|(0;1]$ 与我们的\trBasePath{}$\chi$ 不相交；从现在起，凡考虑 $\pi$ 的\trVanishingPath{}时，这两条都作为常设假定。设 $\widetilde{\Delta}_\gamma$ 是相伴的\trLefschetzThimble{}$\Delta_\gamma$ 在 $\widetilde{E}$ 中的\trPreimage。我们为其配备一个任意的\trGrading{}$\tilde{\alpha}^{\#}$，以及（在\trIsomorphic{}的意义下）唯一的\trPinStructure{}$\widetilde{P}^{\#}$。

要把我们的膜做成\trEquivariant{}的，剩下的事情就是选取一个\trIsomorphism{}$f\colon\widetilde{P}^{\#}\to\iota_{\widetilde{E}}^*\bigl(\widetilde{P}^{\#}\otimes\lambda^{\mathrm{top}}(T\widetilde{\Delta}_\gamma)\bigr)$。按\cref{thm:codim-one-iso-coorientations}，这类\trIsomorphism{}\trCanonical{}地对应于\trFixedPoint{}集的诸\trCoorientation，而在我们的情形，该\trFixedPoint{}集就是 $V_\gamma\subset\widetilde{\Delta}_\gamma$。由于 $\widetilde{\Delta}_\gamma$ \trFibration{}于 $\gamma$ 的\trPreimage{}$\tilde{\gamma}$ 之上，我们有
\begin{equation}\label{eq:dtilde-gamma-normal-line}
D\tilde{\pi}(T\widetilde{\Delta}_\gamma|V_\gamma)=\RR\tilde{\gamma}'(0)=\RR\sqrt{\gamma'(0)},
\end{equation}
从而诸\trCoorientation{}自然地对应于平方根的种种选取。我们的约定是：总取\emph{具有正虚部}的平方根。\eqref{eq:dtilde-gamma-normal-line} 还表明 $\widetilde{\Delta}_\gamma$ 满足\cref{thm:normal-angle-assumption}。最后，$\scr{F}(\widetilde{E})^{\mathrm{naive}}$ 的定义要求对每个 $\widetilde{\Delta}_\gamma$ 选取一个实值\trNormalAngle，我们把它取为 $\arg(\gamma'(0))/2\in(0;\pi)$。

如\cref{sub:grading-and-signs}中所指出的，$\widetilde{\Delta}_\gamma$ 上的一个\trEquivariant{}\trBrane{}结构会在其\trFixedPoint{}集——亦即 $V_\gamma$——上诱导出一个\trBraneStructure。具体地说，该\trGrading{}为
\begin{equation*}
\alpha^{\#}_{V_\gamma}(x)=\tilde{\alpha}^{\#}_{\widetilde{\Delta}_\gamma}(0,x)-\tfrac{1}{2\pi}\arg(\gamma'(0)),
\end{equation*}
而\trPinStructure{}则是（在\trIsomorphic{}的意义下）唯一一个在 $\widetilde{\Delta}_\gamma$ 上界定一个\trPinStructure{}的。这意味着，当 $n=1$ 时，该\trVanishingCycle{}将携带非平凡的\trPinStructure，而不是平凡的那一个（对任何\trCompact\trConnected{}的\trSurface{}$S$，$H^2(S,\partial S;\ZZ/2)\cong\ZZ/2$ 中的相对第二\trStiefelWhitney{}类就是 mod 2 的\trEulerCharacteristic；由于它在 $S=D$ 的情形非零，因此在边界上平凡的\trPinStructure{}是不存在的）。
\end{itemize}
\begin{lemma}\label{thm:type-b-equivariant-hf}
(i) 对任何 (B) 型\trEquivariant{}膜 $\widetilde{\Delta}^{\#}=\widetilde{\Delta}^{\#}_\gamma$，$HF^*(\widetilde{\Delta}^{\#},\widetilde{\Delta}^{\#})\cong H^*(S^{n+1};\KK)$。$\ZZ/2$-作用在 $HF^0$ 上是 $\mathrm{Id}$，而在 $HF^{n+1}$ 上是 $-\mathrm{Id}$。
(ii) 取一对\trVanishingPath{}$\gamma_k$（$k=0,1$），并设 $\widetilde{\Delta}^{\#}_k$ 是与之相伴的 (B) 型\trEquivariant{}膜。那么，\trPoincareDuality{}映射 $HF^*(\widetilde{\Delta}^{\#}_1,\widetilde{\Delta}^{\#}_0)\cong HF^{n+1-*}(\widetilde{\Delta}^{\#}_0,\widetilde{\Delta}^{\#}_1)^{\vee}$ 与诸 $\ZZ/2$-作用反\trCompatible（也就是说，它在相差一个整体 $-1$ 正负号的意义下是\trEquivariant{}的）。
\end{lemma}
\begin{proof}
这些纯粹是\trCohomology{}层面的陈述，可以由\cref{sub:equivariant-donaldson-fukaya}中所列的\trEquivariant{}\trFloerCohomology{}的一般性质推出。事实上，$\iota_{\widetilde{E}}|\widetilde{\Delta}$ 是 $\widetilde{\Delta}$ 的一个\trOrientationReversing{}\trAutomorphism，而 PSS~\trIsomorphism{}与 $\ZZ/2$-作用\trCompatible，故 (i) 成立。至于 (ii)，我们回忆\trPoincareDuality{}\trIsomorphism{}来自\trTriangleProduct
\begin{equation*}
HF^*(\widetilde{\Delta}_1,\widetilde{\Delta}_0)\otimes HF^{n+1-*}(\widetilde{\Delta}_0,\widetilde{\Delta}_1)
\longrightarrow HF^{n+1}(\widetilde{\Delta}_0,\widetilde{\Delta}_0)\cong H^{n+1}(\widetilde{\Delta}_0;\KK)\cong\KK,
\end{equation*}
因此更\trCanonical{}的写法应当是 $HF^*(\widetilde{\Delta}_1,\widetilde{\Delta}_0)\cong H^{n+1}(\widetilde{\Delta}_0;\KK)\otimes HF^{n+1-*}(\widetilde{\Delta}_0,\widetilde{\Delta}_1)^{\vee}$。在这种形式下，它与诸 $\ZZ/2$-作用同样是\trCompatible{}的；而去掉 $H^{n+1}(\widetilde{\Delta}_0;\KK)$ 这一因子则会引进一个正负号，理由正是 (i)。
\end{proof}
一如 (U) 型膜的情形，\trFloerDatum{}与\trPerturbationDatum{}的特种选取有助于让 $\ZZ/2$-作用在上链层面上变得更透明。取两条\trVanishingPath{}$\gamma_k$，以及与之相伴的 (B) 型膜 $\widetilde{\Delta}^{\#}_k$ 和\trVanishingCycle{}$V^n_k$。设
\begin{equation}\label{eq:two-vanishing-paths-position}
\left\{
\begin{aligned}
&\arg(\gamma_0'(0))>\arg(\gamma_1'(0)),\\
&\gamma_0([0;1])\cap\gamma_1([0;1])=\{*\},\ \text{以及}\\
&V_0\pitchfork V_1.
\end{aligned}
\right.
\end{equation}
在为 $(\widetilde{\Delta}_0,\widetilde{\Delta}_1)$ 选取\trFloerDatum{}$(H,J)$ 时，我们想施加一个大大强化了的\cref{thm:floer-datum-assumption}版本，即：每个 $J_t$ 都应使 $\tilde{\pi}$ 成为 $J_t$-\trHolomorphic{}的，并且 $H=0$。后一部分与\cref{thm:floer-datum-assumption}\trCompatible，因为由 \eqref{eq:two-vanishing-paths-position}，诸\trNormalAngle{}满足 $\sigma_1-\sigma_0\in(-\pi;0)$，从而可以在 \eqref{eq:normal-angle-chamber} 中取 $h=0$。在 $\tilde{\pi}$ 之下，任何交点 $\tilde{y}=(y,x)\in\scr{C}(\widetilde{\Delta}_0,\widetilde{\Delta}_1)$ 都必然\trProjection{}到 $\tilde{\gamma}_0$ 与 $\tilde{\gamma}_1$ 的唯一交点。于是 $x=0$，这意味着 $\tilde{y}$ 落在 $\widetilde{E}^{\mathrm{fix}}=M$ 之内。按定义，每个 $\widetilde{\Delta}_k$ 上的\trEquivariant{}\trPinStructure{}对应于 $\widetilde{\Delta}^{\mathrm{fix}}_k$ 的\trNormalBundle{}的这样一个\trTrivialization：它\trProjection{}到具有正虚部的平方根 $\sqrt{\gamma'_k(0)}$。这个平方根是 $\exp(i\sigma_k)=\exp(i\,\arg(\gamma'_k(0))/2)$ 的一个正实数倍。因此，\cref{thm:epsilon-invariant-part}中的诸正负号 $w_k$ 都是 $+1$，这意味着 $\ZZ/2$-作用在 $|o(\tilde{y})|_\KK$ 上是平凡的。这表明，整个\trFloer{}链群 $CF^*(\widetilde{\Delta}^{\#}_0,\widetilde{\Delta}^{\#}_1)$ 都是 $\ZZ/2$-不变的。
类似地，任何 $\tilde{u}\in\scr{M}_Z(\tilde{y}_0,\tilde{y}_1)$ 都\trProjection{}为一张\trHolomorphic{}映射 $\tilde{v}=\tilde{\pi}\circ\tilde{u}\colon Z\to\widetilde{S}$，满足 $\tilde{v}(\RR\times\{k\})\subset\tilde{\gamma}_k([-1;1])$、$\lim_{s\to\pm\infty}\tilde{v}(s,\cdot)=0$。任何这种 $\tilde{v}$ 的\trEnergy{}都必然为零，于是它必是常值的，这意味着 $\tilde{u}$ 同样含于 $M$ 之内。假定 $(H,J)$ 的选取使得 $M$ 上的诸\trFloer{}\trModuliSpace{}都是\trRegular{}的，那么\cref{thm:regularity-fixed-set}便蕴含 $\widetilde{E}$ 中的\trRegularity。此外，\cref{thm:equivariant-index-orientation}表明，无论我们在 $M$ 中还是在 $\widetilde{E}$ 中工作，绝对\trIndex{}与\trOrientation{}空间都是一样的，而 \eqref{eq:equivariant-cu-diagram} 则就每个\trIsolated{}的 $\tilde{u}$ 对\trFloer\trDifferential{}作出贡献时的正负号说了同样的话。总而言之，我们有诸\trComplex{}之间的一个\trIsomorphism
\begin{equation}\label{eq:cf-delta-vanishing-iso}
CF^*(\widetilde{\Delta}^{\#}_0,\widetilde{\Delta}^{\#}_1)\cong CF^*(V^{\#}_0,V^{\#}_1),
\end{equation}
其中右端使用的是从诸 $\widetilde{\Delta}^{\#}_k$ 诱导的、诸\trVanishingCycle{}上的\trBraneStructure。此外，左端上的\trCanonical{}$\ZZ/2$-作用是平凡的。更一般地，设有\trVanishingPath{}$\gamma_0,\ldots,\gamma_d$（$d\ge 2$）使得
\begin{equation}\label{eq:d-vanishing-paths-position}
\left\{
\begin{aligned}
&\arg(\gamma'_k(0))>\arg(\gamma'_{k+1}(0))\ \text{对一切 }k,\\
&\gamma_j([0;1])\cap\gamma_k([0;1])=\{*\}\ \text{对 }j\neq k,\\
&\text{且诸 }V_k\text{ 处于\trGeneralPosition{}}.
\end{aligned}
\right.
\end{equation}
然后，选好满足 $K=0$ 且 $\tilde{\pi}$ 为 $J$-\trHolomorphic{}的\trPerturbationDatum{}$(K,J)$ 之后，便会发现所有\trHolomorphic{} $(d{+}1)$ 边形都留在 $M$ 之内。假定\trRegularity{}成立（见\cref{thm:regularity-polygons}），它们同样会以相同的方式对 $M$ 与 $\widetilde{E}$ 上的\trAinfty-结构作出贡献。
\begin{remark}\label{thm:dual-action-minus-id}
鉴于\cref{thm:type-b-equivariant-hf}(ii)，\eqref{eq:cf-delta-vanishing-iso} 蕴含：对偶群 $HF^*(\widetilde{\Delta}^{\#}_1,\widetilde{\Delta}^{\#}_0)$ 上的\trCanonical{}$\ZZ/2$-作用是 $-\mathrm{Id}$。原则上，人们可以在上链层面证明这一点，办法是取一个\trInhomogeneous{}项为零的\trFloerDatum。而为使它满足 \eqref{eq:normal-angle-chamber}（其中必然 $h=0$），就得把\trNormalAngle{}$\sigma_1$ 换成 $\sigma_1+\pi$，而这又意味着在\cref{thm:epsilon-invariant-part}中 $w_1=-1$，从而使相关\trOrientation{}空间上的 $\ZZ/2$-作用变得非平凡。遗憾的是，\trNormalAngle{}的这一改动使得它与先前对 $(\widetilde{\Delta}_0,\widetilde{\Delta}_1)$ 所用的\trFloerDatum{}不\trCompatible{}了。具体地说，想同时使用二者，就会在\trPseudoholomorphic\trPolygon{}的\trEquivariant{}\trTransversality{}方面遇到潜在的问题，因此我们将放弃这一做法。这意味着，在我们版本的 $\scr{F}(\widetilde{E})^{\mathrm{naive}}$ 中，上链层面的等式 \eqref{eq:cf-delta-vanishing-iso} 成立，而它在 $(\widetilde{\Delta}_1,\widetilde{\Delta}_0)$ 情形的对应物则仅在\trCohomology{}层面上为真。这一选择是由下文的\cref{thm:fibration-exceptional-collection}所决定的。
\end{remark}
\begin{remark}\label{thm:mixed-floer-cohomology}
与 \eqref{eq:tilde-cf-decomposition} 类似的一个结果在混合情形下成立，即其中一个\trSubmanifold{}$\widetilde{L}=\beta_{\widetilde{E}}^{-1}(L)$ 是 (U) 型的，而另一个 $\widetilde{\Delta}=\beta_{\widetilde{E}}^{-1}(\Delta)$ 是 (B) 型的。首先，尽管 $\Delta$ 带边界，$(L,\widetilde{\Delta})$ 的\trFloerCohomology{}仍可按下述方式轻易定义。取一个在 $E^{\mathrm{shrink}}$ 之外平凡的\trFloerDatum{}$(H,J)$，于是对任何 $u\in\scr{M}_Z(y_0,y_1)$，$v=\pi\circ u$ 在 $v^{-1}(S^{\mathrm{shrink}})$ 之外是\trHolomorphic{}的。此时，$v$ 的像不含 $*$；因为若不然，由开性，它必定还含有邻近的点 $z\notin\gamma([0;1])$，这意味着\trIntersectionNumber{}$u\cdot[E_z]$ 将会是正的。由于 $S\setminus\bigl(S^{\mathrm{shrink}}\cup v([0;1])\bigr)$ 是\trConnected{}的，同样的正性对 $z\in\partial S$ 也会成立，这就矛盾了。归结起来：由于 $u$ 不与 $\partial\Delta$ 相交，这一边界的存在性在\trFloer{}理论的考量中可以忽略不计。记住了这一点，便可利用\trIsotopy{}\trInvariance{}把我们的问题化归到 $\gamma(t)=*-t$（$t\in[0;\delta]$）而 $\gamma$ 的其余部分留在 $S^{\mathrm{shrink}}$ 之内的情形。在这种情形下，\trMaximumPrinciple{}照旧适用于 $|v|$，从而给出
\begin{equation*}
HF^*(\widetilde{L},\widetilde{\Delta})\cong HF^*(L,\Delta)\otimes_\KK\KK[+,-].
\end{equation*}
在这一等式的右端，$\Delta$ 带有（在\trIsomorphic{}意义下）唯一的那个平凡\trPinStructure，以及从 $\widetilde{\Delta}$ 的\trGrading{}下降而来的\trGrading，这恰好是把 \eqref{eq:tilde-l-grading} 的过程反过来。
\end{remark}
